
% Copyright 2022 by Robert Hildebrand
%This work is licensed under a
%Creative Commons Attribution-ShareAlike 4.0 International License (CC BY-SA 4.0)
%See http://creativecommons.org/licenses/by-sa/4.0/

\chapter{Duality}
\todoChapter{ {\color{gray}0\% complete. Goal 80\% completion date: January 20, 2023}\\
Notes: This is a borrowed section.  Likely we should update this to match out CC-BY-SA 4.0 license.  Also, update all content to match notation in the book.}
\begin{outcome}
You will learn
\begin{itemize}
\item How to derive a bound on the optimal objective value
\item How to reformulate the problem into a \emph{dual} with a different interpretation
\item Correspondence between the primal (original) and dual problems
\end{itemize}
\end{outcome}

\section{Duality Theory}

\subsection*{Motivating Example}

Imagine you run a small bakery that makes two products: cakes ($x_1$) and cookies ($x_2$). Your goal is to maximize your daily profit. Cakes yield a profit of \$3 each, and cookies yield \$2 each. You are limited by the ingredients you have on hand: you have at most 4 units of flour and 5 units of sugar available each day. Each cake requires 1 unit of flour and 2 units of sugar, while each cookie requires 1 unit of flour and 1 unit of sugar. We can write this scenario as a linear program:

\[
\begin{aligned}
\max \; & 3x_1 + 2x_2 && \text{(maximize total profit)}\\
\text{subject to } & x_1 + x_2 \leq 4 && \text{(flour constraint)}\\
& 2x_1 + x_2 \leq 5 && \text{(sugar constraint)}\\
& x_1, x_2 \geq 0.
\end{aligned}
\]

Here, $x_1$ is the number of cakes you bake, and $x_2$ is the number of cookies. The constraints limit how many of each product you can produce given your resources.

Now, suppose a friend is not sure that you are telling the truth about the maximum profit you can make. They ask, "Can you prove an upper bound on how much money you could possibly earn with these ingredient constraints, without actually solving the problem?"

\subsection*{Intuition Behind Duality}

Consider the constraints in the problem. Each one puts a cap on how large the combination of $x_1$ and $x_2$ can be. If we think of them separately:

\begin{itemize}
\item The flour constraint: $x_1 + x_2 \leq 4$.
\item The sugar constraint: $2x_1 + x_2 \leq 5$.
\end{itemize}
Now, what if we take some nonnegative multipliers $y_1$ and $y_2$ and combine these inequalities into one single inequality? For instance, multiply the flour constraint by $y_1 \geq 0$ and the sugar constraint by $y_2 \geq 0$ and add them up:

\[
y_1(x_1 + x_2) + y_2(2x_1 + x_2) \leq y_1(4) + y_2(5).
\]

Grouping the terms by $x_1$ and $x_2$ gives:

\[
(x_1(y_1 + 2y_2)) + (x_2(y_1 + y_2)) \leq 4y_1 + 5y_2.
\]

If we choose $y_1$ and $y_2$ cleverly, we can force the left-hand side to be an expression that is always \emph{at least} as large as our profit function $3x_1 + 2x_2$. Specifically, we want

\[
y_1 + 2y_2 \geq 3 \quad \text{and} \quad y_1 + y_2 \geq 2.
\]

Why? Because if the coefficients on $x_1$ and $x_2$ in our combined inequality dominate those in our profit function, then the upper bound on the combined inequality's right-hand side ($4y_1 + 5y_2$) will also serve as an upper bound on the profit $3x_1 + 2x_2$. This ensures:

\[
3x_1 + 2x_2 \leq (y_1+2y_2)x_1 + (y_1 + y_2)x_2 \leq 4y_1 + 5y_2.
\]

Thus, by choosing $y_1$ and $y_2$, we obtain a valid upper bound on the maximum profit. If we try to minimize $4y_1 + 5y_2$ over all $y_1, y_2 \geq 0$ that satisfy these inequalities, we find the best (smallest) upper bound on our original problem.

\subsection*{General Derivation of the Dual}

We started with the primal problem in a general form:

\[
\begin{aligned}
\max \; & \c^\top \x \\
\text{subject to } & A\x \leq \b \\
& \x \geq 0.
\end{aligned}
\]

Each constraint $a^i \cdot x \leq b_i$ (where $a^i$ is the $i$-th row of $A$) is valid for all feasible $x$. If we take a nonnegative combination of these constraints, say with multipliers $y_i \geq 0$, we get:

\[
\sum_{i=1}^m y_i (a^i \cdot \x) \leq \sum_{i=1}^m y_i b_i.
\]

This can be written as:

\[
(\sum_{i=1}^m y_i \a^i) \cdot \x \leq \sum_{i=1}^m y_i b_i.
\]

If we choose $\y$ such that $\sum_{i=1}^m y_i a^i \geq c$, then for all feasible $\x$ (since $\x \geq 0$):

\[
\c^\top \x \leq (\sum_{i=1}^m y_i a^i)\cdot \x \leq \sum_{i=1}^m y_i b_i.
\]

So, $\b^\top \y = \sum_{i=1}^m y_i b_i$ is an upper bound on the maximum primal objective value. By minimizing over all such $y$, we seek the tightest upper bound:

\[
\begin{aligned}
\min \; & \b^\top \y \\
\text{subject to } & A^\top \y \geq \c \\
& \y \geq 0.
\end{aligned}
\]

This is the \emph{dual} problem. In our bakery scenario, this dual formulation corresponds to adjusting the prices (the $y$ values can be thought of as "shadow prices") you would assign to each ingredient such that if you charge yourself for using these ingredients at these shadow prices, you end up with the smallest possible "bill" that still always surpasses or equals your maximum possible profit. The beauty of the Strong Duality Theorem ensures that the optimal primal value equals the optimal dual value.


From the above reasoning, we have shown that if $x$ is feasible for the primal and $y$ is feasible for the dual, then $c^\top x \leq b^\top y$. In other words, the value of any feasible primal solution is always bounded above by the value of any feasible dual solution.

\begin{theorem}{Weak Duality}{weak-duality}
For any feasible $\x$ in the primal problem and any feasible $\y$ in the dual problem, we have:
\[
\c^\top \x \leq \b^\top \y.
\]
\end{theorem}

The weak duality theorem is not only specific to linear programming problems; it appears in many other classes of optimization problems. It establishes an important inequality between primal and dual objective values for any pair of feasible solutions.

However, in the special case of linear programs (and certain other problem classes), we can say even more. In these settings, not only does the dual objective always provide an upper bound on the primal objective, but there exist optimal solutions for both problems whose objective values match. This result is known as the \emph{Strong Duality Theorem}.

\section{Primal-Dual Pairs and Strong Duality}
\begin{outcome}
\begin{itemize}
\item Learn the two most important theorems in Linear Programming!!!
\end{itemize}
\end{outcome}
Consider the following pair of linear programs, known as a primal-dual pair. The primal problem (P) has $n$ decision variables and $m$ constraints, while the dual problem (D) has $m$ decision variables and $n$ constraints.
\begin{multicols}{2}
\begin{tcolorbox}[title=Primal (P), colback=blue!5, colframe=blue!60!black, sharp corners=south]
\[
\begin{aligned}
\max\ & \mathbf{c}^\top \mathbf{x} \\
\text{s.t.} \quad & A \mathbf{x} \leq \mathbf{b} \\
& \mathbf{x} \geq 0 \\
& \text{(Here, } \mathbf{x} \in \mathbb{R}^n \text{ and } A \text{ is an } m \times n \text{ matrix.)}
\\
\phantom{} \\
\end{aligned}
\]
\end{tcolorbox}

\columnbreak

\begin{tcolorbox}[title=Dual (D), colback=green!5, colframe=green!60!black, sharp corners=south]
\[
\begin{aligned}
\min\ & \mathbf{b}^\top \mathbf{y} \\
\text{s.t.} \quad & A^\top \mathbf{y} \geq \mathbf{c} \\
& \mathbf{y} \geq 0 \\
& \text{(Here, } \mathbf{y} \in \mathbb{R}^m, \text{ so the dual has } \\
& m \text{ variables and } n \text{ constraints.)}
\end{aligned}
\]
\end{tcolorbox}
\end{multicols}

From weak duality, we know that for any feasible $\x$ in (P) and any feasible $\y$ in (D), it holds that:
\[
\c^\top \x \leq \b^\top \y.
\]

This indicates that the optimal value of the dual problem is always an upper bound on the optimal value of the primal problem.

\begin{theorem}{Weak Duality}{weak-duality}
For any feasible solution $\x$ of the primal problem (P) and any feasible solution $\y$ of the dual problem (D), 
\[
\c^\top \x \leq \b^\top \y.
\]
\end{theorem}

Weak duality theorems appear in many classes of optimization problems. In the special case of linear programs (and some other structured problems), we also have a stronger result, known as the \emph{Strong Duality Theorem}. This theorem completely characterizes the relationship between the primal and dual solutions:

\begin{theorem}{Strong Duality}{strong-duality}
For a primal-dual pair of linear programs:
\[
\begin{aligned}
\textbf{(P)} \quad &\max \c^\top \x \\
&\text{subject to } A \x \leq \b, \x \geq 0,
\end{aligned}
\quad\quad
\begin{aligned}
\textbf{(D)} \quad &\min \b^\top \y \\
&\text{subject to } A^\top \y \geq \c, \y \geq 0,
\end{aligned}
\]
the following cases hold:
\begin{enumerate}
\item  If (P) is unbounded, then (D) must be infeasible.
\item  If (D) is unbounded, then (P) must be infeasible.
\item  If both (P) and (D) have feasible solutions and are bounded, then both have optimal solutions $\x^*$ and $\y^*$ respectively, and 
\[
\c^\top \x^* = \b^\top \y^*.
\]
\end{enumerate}
\end{theorem}

\begin{example}{Dual Infeasibility from Primal Unboundedness}{}\label{example:simplex-unbounded}
Consider the following linear program:
\[
\begin{aligned}
\text{maximize} \quad & x_1 \\
\text{subject to} \quad
& x_1 - x_2 \leq 1 \\
& 2x_1 - x_2 \leq 3 \\
& x_1, x_2 \geq 0
\end{aligned}
\]

As shown in Example~\ref{example:simplex-unbounded}, this primal problem is unbounded. After applying the simplex method, we found that when \( s_1 \) enters the basis, there is no constraint preventing it from growing indefinitely, and the objective increases without bound. Therefore, the primal is unbounded.

Now, consider the dual of this problem:
\[
\begin{aligned}
\text{minimize} \quad & y_1 + 3y_2 \\
\text{subject to} \quad
& y_1 + 2y_2 \geq 1 \\
& -y_1 - y_2 \geq 0 \\
& y_1, y_2 \geq 0
\end{aligned}
\]
\textbf{Conclusion via Strong Duality:}  
Because the primal problem is unbounded and the primal-dual pair satisfies the usual assumptions of linear programming, \textbf{strong duality} implies that the dual must be \textbf{infeasible}. 


This is confirmed by the Big-M simplex pivot steps, which fail to drive the artificial variable to zero.
To solve the dual, we apply the simplex method using the Big-M technique to initialize feasibility. The auxiliary variable \( a_1 \) is introduced to handle the infeasibility of the initial basic solution.

\textbf{First dictionary (with Big-M penalty):}
\[
\text{Maximize } -y_1 - 3y_2 - 1000a_1
\]
\[
\text{subject to} \quad
\begin{cases}
w_1 = -1 + y_1 + 2y_2 + a_1 \\
w_2 = 0 - y_1 - y_2
\end{cases}
\]

\textbf{Pivot 1:} Enter \( y_1 \), leave \( a_1 \). We rewrite in terms of the new basic variables.

\textbf{Second dictionary:}
\[
\text{Maximize } -1000 + 999y_1 + 1997y_2 - 1000w_1
\]
\[
\text{subject to} \quad
\begin{cases}
a_1 = 1 - y_1 - 2y_2 + w_1 \\
w_2 = 0 - y_1 - y_2
\end{cases}
\]

\textbf{Pivot 2:} Enter \( y_2 \), leave \( w_2 \). Continue pivoting.

\textbf{Final dictionary:}
\[
\text{Maximize } -1000 - 998y_1 - 1000w_1 - 1997w_2
\]
\[
\text{subject to} \quad
\begin{cases}
a_1 = 1 + y_1 + w_1 + 2w_2 \\
y_2 = 0 - y_1 - w_2
\end{cases}
\]

In this final dictionary, all variables in the objective function have non-positive coefficients, but the artificial variable \( a_1 \) remains in the basis with a positive value. Since \( a_1 > 0 \) at termination, we conclude that the original dual problem is \emph{infeasible}.



\end{example}


\subsection{Economic Interpretation of the Dual}

Duality in linear programming offers a powerful economic interpretation: the dual variables represent the \emph{value} of resources or constraints in the primal problem. These values are commonly known as \textbf{shadow prices}.

\vspace{1em}
\begin{definition}{Shadow Price}{}
The \emph{shadow price} of a constraint in a linear program tells us how much the optimal objective value would change if we slightly increased the right-hand side of that constraint by one unit, while keeping everything else fixed.

\smallskip
In a maximization problem, a shadow price represents the \textbf{increase in value} per additional unit of a scarce resource. In a minimization problem, it represents the \textbf{decrease in cost}.

\smallskip
\textbf{Interpretation:}  
If a constraint represents the availability of a resource (e.g., time, material), then its shadow price reflects the marginal value of that resource — how much the objective would improve if you had one more unit of it.

\smallskip
\textbf{Mathematically:}  
If the constraint \( a_i^\top \x \leq b_i \) has corresponding dual variable \( y_i^* \), then:
\[
\text{Shadow price} = \frac{d z^*}{d b_i} = y_i^*
\]
where \( z^* \) is the optimal value of the objective function.
\end{definition}

\paragraph{Units and Interpretation:}
The units of a shadow price match the units of the objective function (e.g., dollars of profit) per unit of the resource (e.g., kilograms of flour, hours of labor). If $y_2 = 10$ in the bakery example, and the constraint corresponds to sugar (in kg), then each additional kg of sugar would allow the bakery to increase profit by \$10.

\vspace{1em}
\noindent\textbf{Example: The Bakery Revisited}

Consider the primal problem:

$$
\begin{aligned}
\text{maximize} \quad & 3x_1 + 2x_2 \\
\text{subject to} \quad
& x_1 + x_2 \leq 4 \quad \text{(Flour)} \\
& 2x_1 + x_2 \leq 5 \quad \text{(Sugar)} \\
& x_1, x_2 \geq 0
\end{aligned}
$$

The dual problem is:

$$
\begin{aligned}
\text{minimize} \quad & 4y_1 + 5y_2 \\
\text{subject to} \quad
& y_1 + 2y_2 \geq 3 \\
& y_1 + y_2 \geq 2 \\
& y_1, y_2 \geq 0
\end{aligned}
$$

If the optimal dual solution is $y_1^* = 1, y_2^* = 1$, then:

\begin{itemize}
\item The shadow price of flour is 1 — each extra unit of flour increases profit by \$1.
\item The shadow price of sugar is 1 — each extra unit of sugar increases profit by \$1.
\end{itemize}

\noindent \textbf{Example: Bakery 2.0}

The primal problem describes a bakery optimizing the production of cakes, cookies, and muffins to maximize profit. The constraints represent limited resources: flour, sugar, and baking time. The dual problem attaches a shadow price to each of these resources, reflecting their marginal value.

\vspace{1em}
\noindent\textbf{Primal Problem:}
$$
\begin{aligned}
\max\quad & 60x_1 + 30x_2 + 20x_3 \\
\text{s.t.} \quad
& 8x_1 + 6x_2 + x_3 \leq 48 \quad \text{(Flour)} \\
& 4x_1 + 2x_2 + 1.5x_3 \leq 20 \quad \text{(Sugar)} \\
& 2x_1 + 1.5x_2 + 0.5x_3 \leq 8 \quad \text{(Baking Time)} \\
& x_1, x_2, x_3 \geq 0
\end{aligned}
$$
\textbf{Optimal Primal Solution:}
$$
x_1 = 2,\quad x_2 = 0,\quad x_3 = 8,\quad \text{Objective Value } = 280
$$
\vspace{1em}
\noindent\textbf{Dual Problem:}
$$
\begin{aligned}
\min\quad & 48y_1 + 20y_2 + 8y_3 \\
\text{s.t.} \quad
& 8y_1 + 4y_2 + 2y_3 \geq 60 \quad \text{(Cakes)} \\
& 6y_1 + 2y_2 + 1.5y_3 \geq 30 \quad \text{(Cookies)} \\
& y_1 + 1.5y_2 + 0.5y_3 \geq 20 \quad \text{(Muffins)} \\
& y_1, y_2, y_3 \geq 0
\end{aligned}
$$
\textbf{Optimal Dual Solution:}
$$
y_1 = 0,\quad y_2 = 10,\quad y_3 = 10,\quad \text{Objective Value } = 280
$$

\vspace{1em}
\subsubsection*{Interpretation of Shadow Prices}
\begin{itemize}
\item $y_1 = 0$: The shadow price of flour is \$0.
This implies the flour constraint is \textbf{non-binding} — the bakery is not using all 48 kg of flour. Increasing flour supply would not increase profit.

\item $y_2 = 10$: The shadow price of sugar is \$10.
Each additional kilogram of sugar allows the bakery to increase profit by \$10. The sugar constraint is \textbf{binding}.

\item $y_3 = 10$: The shadow price of baking time is \$10 per hour.
Baking time is fully used and limits production. Gaining one more hour would allow an increase in profit by \$10.
\end{itemize}

\vspace{1em}
\subsubsection*{Complementary Slackness Verification}

From the primal solution:
\begin{itemize}
\item The \textbf{flour constraint} is not tight \ $\Rightarrow$ \  $y_1 = 0$
\item The \textbf{sugar constraint} is tight \ $\Rightarrow$ \ $y_2 > 0$
\item The \textbf{baking time constraint} is tight \ $\Rightarrow$ \ $y_3 > 0$
\end{itemize}

These relationships satisfy the \emph{complementary slackness conditions}:

$$
y_i > 0 \Rightarrow \text{Constraint } i \text{ is tight}, \quad
\text{Constraint } i \text{ not tight} \Rightarrow y_i = 0
$$

\vspace{1em}
\textbf{Sensitivity Analysis: Marginal Value of Resources}
Assuming the current basis remains optimal, the increase in maximum profit from one additional unit of:
\begin{itemize}
\item Flour = \$0
\item Sugar = \$10
\item Baking Time = \$10
\end{itemize}


\vspace{1em}


The dual solution not only validates the primal objective value (via strong duality), but also assigns economic meaning to each constraint. Shadow prices inform how much you’d be willing to pay for more of each resource — a foundational concept in operations, pricing, and planning.


\paragraph{When Are Shadow Prices Meaningful?}
Shadow prices are meaningful only when:

\begin{itemize}
\item The constraint is \emph{binding} (i.e., tight) at the optimal solution.
\item The problem satisfies strong duality (which always holds for feasible and bounded linear programs).
\item The right-hand side is perturbed slightly (local sensitivity).
\end{itemize}

If a constraint is not tight, then its shadow price is 0 — the resource is not limiting the optimal solution.

\vspace{1em}
\paragraph{Connection to the Dual Problem:}
In the dual problem:
\begin{itemize}
\item The \emph{objective function} is the total value of resources, priced at shadow prices.
\item The \emph{constraints} ensure that no activity in the primal makes more profit than it costs, when inputs are valued at shadow prices.
\end{itemize}
This ensures equilibrium: no profitable arbitrage is possible, and every unit of constrained resource is priced at its true marginal value.

\vspace{1em}
\noindent\textbf{Summary of Shadow Price Properties:}
\begin{itemize}
\item A shadow price represents how much you would \emph{pay} for an extra unit of a constrained resource.
\item It is equal to the dual variable corresponding to the constraint.
\item It reflects the marginal value of relaxing that constraint.
\item Shadow prices are only nonzero for \textbf{binding} constraints.
\end{itemize}

\vspace{0.5em}

\begin{itemize}
\item If a resource is fully used and limits optimality, its shadow price is positive.
\item If a resource is unused or abundant, its shadow price is zero.
\end{itemize}




\section{Different Primal–Dual Forms}

The form of a dual problem depends on the primal's objective (maximize or minimize), inequality directions, and sign constraints on variables. The following layout uses paired boxes to clearly compare each primal and its dual.

\subsection*{1. Maximization with $=$ Constraints (Standard Form)}

\begin{multicols}{2}
\begin{tcolorbox}[title=Primal (P), colback=blue!5, colframe=blue!40!black]
\[
\begin{aligned}
\max\ & \mathbf{c}^\top \mathbf{x} \\
\text{s.t.} \quad & A \mathbf{x} = \mathbf{b} \\
& \mathbf{x} \geq 0
\end{aligned}
\]
\end{tcolorbox}

\columnbreak

\begin{tcolorbox}[title=Dual (D), colback=green!5, colframe=green!40!black]
\[
\begin{aligned}
\min\ & \mathbf{b}^\top \mathbf{y} \\
\text{s.t.} \quad & A^\top \mathbf{y} \geq \mathbf{c} \\
& \mathbf{y} \ \ \text{ free}
\end{aligned}
\]
\end{tcolorbox}
\end{multicols}

\subsection*{2. Minimization with $\leq$ Constraints}

\begin{multicols}{2}
\begin{tcolorbox}[title=Primal (P), colback=blue!5, colframe=blue!40!black]
\[
\begin{aligned}
\min\ & \mathbf{b}^\top \mathbf{x} \\
\text{s.t.} \quad & A \mathbf{x} \leq \mathbf{c} \\
& \mathbf{x} \geq 0
\end{aligned}
\]
\end{tcolorbox}

\columnbreak

\begin{tcolorbox}[title=Dual (D), colback=green!5, colframe=green!40!black]
\[
\begin{aligned}
\max\ & \mathbf{c}^\top \mathbf{y} \\
\text{s.t.} \quad & A^\top \mathbf{y} \leq \mathbf{b} \\
& \mathbf{y} \geq 0
\end{aligned}
\]
\end{tcolorbox}
\end{multicols}

\subsection*{3. Equality Constraints in the Primal}

\begin{multicols}{2}
\begin{tcolorbox}[title=Primal (P), colback=blue!5, colframe=blue!40!black]
\[
\begin{aligned}
\max\ & \mathbf{c}^\top \mathbf{x} \\
\text{s.t.} \quad & A \mathbf{x} \leq \mathbf{b} \\
& \mathbf{x} \ \ \text{free}
\end{aligned}
\]
\end{tcolorbox}

\columnbreak

\begin{tcolorbox}[title=Dual (D), colback=green!5, colframe=green!40!black]
\[
\begin{aligned}
\min\ & \mathbf{b}^\top \mathbf{y} \\
\text{s.t.} \quad & A^\top \mathbf{y} = \mathbf{c} \\
& \mathbf{y} \geq 0
\end{aligned}
\]
\end{tcolorbox}
\end{multicols}

\subsection*{4. Sign Variations in Primal Variables and Constraints}

\begin{itemize}
  \item If a primal variable \( x_j \) is unrestricted (free), the corresponding dual constraint is an \emph{equality}.
  \item If \( x_j \leq 0 \), the dual constraint becomes \( \leq \) instead of \( \geq \).
  \item If a primal constraint \( a_i^\top \mathbf{x} \leq b_i \) is changed to \( \geq b_i \), then the sign constraint on the corresponding dual variable \( y_i \) flips.
\end{itemize}

\textbf{Example: Primal with Free Variables and Equalities}

\begin{multicols}{2}
\begin{tcolorbox}[title=Primal (P), colback=blue!5, colframe=blue!40!black]
\[
\begin{aligned}
\max\ & \mathbf{c}^\top \mathbf{x} \\
\text{s.t.} \quad & A \mathbf{x} \leq \mathbf{b} \\
& \mathbf{x} \text{ free}
\end{aligned}
\]
\end{tcolorbox}

\columnbreak

\begin{tcolorbox}[title=Dual (D), colback=green!5, colframe=green!40!black]
\[
\begin{aligned}
\min\ & \mathbf{b}^\top \mathbf{y} \\
\text{s.t.} \quad & A^\top \mathbf{y} = \mathbf{c} \\
& \mathbf{y} \geq 0
\end{aligned}
\]
\end{tcolorbox}
\end{multicols}

\bigskip

This structured approach helps you identify how each component in the primal affects the dual—particularly how inequality direction and sign restrictions control the dual variable domain and constraint form.


As you can see, changing the primal’s constraints and variable domains changes the structure of the dual problem. Every such transformation has a corresponding primal-dual pair that maintains the duality principles. In each case, the weak and strong duality theorems still apply.


\subsection*{Deriving the Dual with Equality Constraints}

Let us now consider a variation of our previous derivation of the dual problem. This time, our primal problem (P) will have equality constraints rather than inequalities, and all variables will remain nonnegative. For example, consider the following linear program with $m=2$ equality constraints and $n=4$ variables:

\[
\begin{aligned}
\max \; & 2x_1 + 3x_2 - x_3 + x_4 \\
\text{subject to } & x_1 + x_2 + 0 \cdot x_3 + x_4 = 10 \\
& 2x_1 + x_2 - x_3 + 0 \cdot x_4 = 5 \\
& x_1, x_2, x_3, x_4 \geq 0.
\end{aligned}
\]

In matrix form, if we let 
\[
A = \begin{bmatrix}
1 & 1 & 0 & 1 \\
2 & 1 & -1 & 0
\end{bmatrix}, \quad b = \begin{bmatrix}10 \\ 5\end{bmatrix}, \quad c = \begin{bmatrix}2 \\ 3 \\ -1 \\ 1\end{bmatrix},
\]
our primal problem is:
\[
\max \c^\top \x \quad\text{subject to } A \x = \b, \; \x \geq 0.
\]

\subsubsection*{Forming a Bound Using Combinations of Equalities}

Previously, when we had inequalities, we formed upper bounds on the objective by taking nonnegative linear combinations of the constraints. Now, since we have \emph{equalities}, we are free to choose any real multipliers for these constraints (positive, negative, or zero), and the resulting combination will still hold as an equality. This is because any scalar multiple of an equation remains a valid equation.

Let $y_1$ and $y_2$ be arbitrary real numbers, and form the linear combination:
\[
y_1 (x_1 + x_2 + x_4 = 10) \quad \text{and} \quad y_2 (2x_1 + x_2 - x_3 = 5).
\]
Adding these together gives:
\[
y_1 x_1 + y_1 x_2 + y_1 x_4 \;+\; y_2(2x_1 + x_2 - x_3) \;=\; 10y_1 + 5y_2.
\]

Collecting terms by variable:
\[
x_1(y_1 + 2y_2) + x_2(y_1 + y_2) + x_3(-y_2) + x_4(y_1) = 10y_1 + 5y_2.
\]

Since $x \geq 0$, if we want the left-hand side to serve as an upper bound on our objective function $2x_1 + 3x_2 - x_3 + x_4$ for all feasible $x$, we need the coefficients of each $x_j$ in the combined inequality to dominate those in the objective. Specifically, we want:
\[
y_1 + 2y_2 \geq 2, \quad y_1 + y_2 \geq 3, \quad -y_2 \geq -1 \;(\text{which implies } y_2 \leq 1), \quad y_1 \geq 1.
\]

If these inequalities hold, then for every feasible $x$:
\[
2x_1 + 3x_2 - x_3 + x_4 \;\leq\; (y_1+2y_2)x_1 + (y_1+y_2)x_2 + (-y_2)x_3 + (y_1)x_4 \;=\; 10y_1 + 5y_2.
\]

Thus, $10y_1 + 5y_2$ is an upper bound on the maximal value of $2x_1 + 3x_2 - x_3 + x_4$. Since $y_1$ and $y_2$ are unrestricted in sign, we can try to \emph{minimize} $10y_1 + 5y_2$ subject to ensuring these dominance conditions.

\subsubsection*{The Dual Problem}

In the general setting with equality constraints and nonnegative primal variables, the dual variables associated with these equality constraints are \emph{free} (they can be positive, negative, or zero). The process we followed above leads us to the dual problem:
\[
\min \b^\top \y = 10y_1 + 5y_2
\]
subject to the condition that $A^\top \y \geq \c$, i.e.:
\[
\begin{bmatrix}
1 & 2 \\
1 & 1 \\
0 & -1 \\
1 & 0
\end{bmatrix}
\begin{bmatrix}y_1 \\ y_2\end{bmatrix} 
\geq
\begin{bmatrix}2 \\ 3 \\ -1 \\ 1\end{bmatrix}.
\]

This system (in our specific numeric example) determines a unique $y$ that perfectly matches the $c$ vector. More generally, the dual problem corresponding to
\[
\max \c^\top \x \quad\text{subject to } A \x = \b, \x \geq 0
\]
is
\[
\min \b^\top \y \quad \text{subject to } A^\top \y \geq \c,\; \y \text{ free}.
\]

The absence of sign restrictions on $\y$ is a direct consequence of using equalities in the primal. Since no direction of inequality is fixed, you are free to "add" or "subtract" the given equations as you please, hence no sign constraints emerge on the dual variables.


\subsection*{Example with Inequalities and Free Variables}

Now consider a primal problem with three inequality constraints and two free variables. Suppose:

\begin{multicols}{2}
\begin{tcolorbox}[title=Primal (P), colback=blue!5, colframe=blue!60!black, sharp corners=south]
\[
\begin{aligned}
\max\ & 3x_1 - 2x_2 \\
\text{s.t.} \quad 
& x_1 + x_2 \leq 10 \\
& 2x_1 - x_2 \leq 5 \\
& -x_1 + 3x_2 \leq 7 \\
& x_1 \text{ free, } x_2 \text{ free}
\end{aligned}
\]
\end{tcolorbox}

\columnbreak

\begin{tcolorbox}[title=Dual (D), colback=green!5, colframe=green!60!black, sharp corners=south]
\[
\begin{aligned}
\min\ & 10y_1 + 5y_2 + 7y_3 \\
\text{s.t.} \quad 
& y_1 + 2y_2 - y_3 = 3 \quad \text{(coeff. of } x_1) \\
& y_1 - y_2 + 3y_3 = -2 \quad \text{(coeff. of } x_2) \\
& y_1, y_2, y_3 \geq 0
\end{aligned}
\]
\end{tcolorbox}
\end{multicols}

Here, $x_1$ and $x_2$ can be positive or negative. 

we can write this in a compact matrix form. Let:
\[
A = \begin{bmatrix}
1 & 1 \\
2 & -1 \\
-1 & 3 
\end{bmatrix}, \quad b = \begin{bmatrix}
10 \\ 5 \\ 7
\end{bmatrix}, \quad c = \begin{bmatrix}
3 \\ -2
\end{bmatrix}.
\]

Then (P) is:
\[
\max \c^\top \x \quad \text{subject to } A \x \leq \b, \text{ and } \x \text{ free}.
\]



To create a bound on $3x_1 - 2x_2$, we take nonnegative combinations of the constraints. Let $y_1, y_2, y_3 \geq 0$ be the multipliers for the three constraints:

\[
y_1(x_1 + x_2) + y_2(2x_1 - x_2) + y_3(-x_1 + 3x_2) \leq 10y_1 + 5y_2 + 7y_3.
\]

Combining like terms in $x_1$ and $x_2$:

\[
x_1(y_1 + 2y_2 - y_3) + x_2(y_1 - y_2 + 3y_3) \leq 10y_1 + 5y_2 + 7y_3.
\]

We want this inequality to hold for all $x_1, x_2$ \emph{free}. If $y_1 + 2y_2 - y_3 > 3$, then by choosing $x_1$ large and positive, we could make the left side arbitrarily large, violating the bound. If $y_1 + 2y_2 - y_3 < 3$, by choosing $x_1$ large and negative, we break the bound. The only way to always hold the inequality is to have exact equality:

\[
y_1 + 2y_2 - y_3 = 3.
\]

Similarly, for $x_2$:

If $y_1 - y_2 + 3y_3 > -2$, we could send $x_2$ to $-\infty$ to break the bound. If $y_1 - y_2 + 3y_3 < -2$, we could send $x_2$ to $\infty$. Again, the only safe choice is equality:

\[
y_1 - y_2 + 3y_3 = -2.
\]

Thus, the dual constraints must be equalities. The dual problem is:

\[
\textbf{Dual (D):} \quad
\begin{aligned}
\min \; & 10y_1 + 5y_2 + 7y_3 \\
\text{subject to } & y_1 + 2y_2 - y_3 = 3 \\
& y_1 - y_2 + 3y_3 = -2 \\
& y_1, y_2, y_3 \geq 0.
\end{aligned}
\]

From our derivation, the dual problem (D) becomes:
\[
\min \b^\top \y = 10y_1 + 5y_2 + 7y_3
\]
subject to:
\[
\begin{bmatrix}
1 & 2 & -1 \\
1 & -1 & 3
\end{bmatrix}
\begin{bmatrix}
y_1 \\ y_2 \\ y_3
\end{bmatrix}
=
\begin{bmatrix}
3 \\ -2
\end{bmatrix},
\]
and
\[
y_1, y_2, y_3 \geq 0.
\]

Summarizing, we have

\begin{multicols}{2}
\begin{tcolorbox}[title=Primal (P), colback=blue!5, colframe=blue!60!black, sharp corners=south]
\[
\begin{aligned}
\max \quad &
\begin{bmatrix} 3 & -2 \end{bmatrix}
\begin{bmatrix} x_1 \\ x_2 \end{bmatrix} \\[0.5em]
\text{s.t.} \quad &
\begin{bmatrix}
1 & 1 \\
2 & -1 \\
-1 & 3
\end{bmatrix}
\begin{bmatrix} x_1 \\ x_2 \end{bmatrix}
\leq
\begin{bmatrix} 10 \\ 5 \\ 7 \end{bmatrix} \\[0.5em]
& x_1, x_2 \text{ free}
\end{aligned}
\]
\end{tcolorbox}

\columnbreak

\begin{tcolorbox}[title=Dual (D), colback=green!5, colframe=green!60!black, sharp corners=south]
\[
\begin{aligned}
\min \quad &
\begin{bmatrix} 10 & 5 & 7 \end{bmatrix}
\begin{bmatrix} y_1 \\ y_2 \\ y_3 \end{bmatrix} \\[0.5em]
\text{s.t.} \quad &
\begin{bmatrix}
1 & 2 & -1 \\
1 & -1 & 3
\end{bmatrix}
\begin{bmatrix} y_1 \\ y_2 \\ y_3 \end{bmatrix}
=
\begin{bmatrix} 3 \\ -2 \end{bmatrix} \\[0.5em]
& y_1, y_2, y_3 \geq 0
\end{aligned}
\]
\end{tcolorbox}
\end{multicols}



Here we see the pattern: having $\x$ free in the primal forces the dual constraints to be equalities. In summary:
\begin{itemize}
\item Primal inequalities with $\x$ free $\Rightarrow$ Dual equality constraints.
\item Primal equalities with $\x \geq 0$ $\Rightarrow$ Dual inequalities and $\y$ free.
\end{itemize}

These relationships ensure that by choosing appropriate dual multipliers, we can always create a valid bound on the primal objective.


The following table can help remember the above.\\

\begin{center}
\begin{tabular}{cc}
\toprule
\textbf{Primal (min)}       & \textbf{Dual (max)}      \\ 
\midrule
\(\geq\) constraint         & \(\geq 0\) variable      \\ 
\(\leq\) constraint         & \(\leq 0\) variable      \\ 
\(=\) constraint            & \(\text{free}\) variable \\ 
\(\geq 0\) variable         & \(\leq\) constraint      \\ 
\(\leq 0\) variable         & \(\geq\) constraint      \\ 
\(\text{free}\) variable    & \(=\) constraint         \\ 
\bottomrule
\end{tabular}
\end{center}

\subsection*{Complicated Dual}
\label{example:complicated-dual}
Consider the following primal linear program:
\begin{multicols}{2}
\begin{tcolorbox}[title=Primal (P), colback=blue!5, colframe=blue!60!black, sharp corners=south]
\[
\begin{aligned}
\max\quad & 2x_1 + x_2 \\
\text{s.t.} \quad
& x_1 + x_2 = 2 \quad \text{(1)} \\
& 2x_1 - x_2 \geq 3 \quad \text{(2)} \\
& x_1 - x_2 \leq 1 \quad \text{(3)} \\
& x_1 \geq 0, \; x_2 \text{ free}
\end{aligned}
\]
\end{tcolorbox}

\columnbreak

\begin{tcolorbox}[title=Dual (D), colback=green!5, colframe=green!60!black, sharp corners=south]
Let \( y_1 \in \mathbb{R} \), \( y_2 \geq 0 \), \( y_3 \leq 0 \) be dual variables for constraints (1), (2), and (3), respectively.  
\[
\begin{aligned}
\min\quad & 2y_1 + 3y_2 + 1y_3 \\
\text{s.t.} \quad
& y_1 + 2y_2 + y_3 \geq 2 \quad \text{(from } x_1 \geq 0\text{)} \\
& y_1 - y_2 - y_3 = 1 \quad \text{(from } x_2 \text{ free)}\\
& y_1  \text{ free } , y_2 \geq 0 , y_3 \leq 0 
\end{aligned}
\]
\end{tcolorbox}
\end{multicols}

To derive the dual, we introduce a dual variable for each constraint:
\begin{itemize}
    \item Let \( y_1 \) be the dual variable for the equality constraint \( x_1 + x_2 = 2 \)
    \item Let \( y_2 \) be the dual variable for the \( \geq \) constraint \( 2x_1 - x_2 \geq 3 \)
    \item Let \( y_3 \) be the dual variable for the \( \leq \) constraint \( x_1 - x_2 \leq 1 \)
\end{itemize}

\textbf{Explanation of Dual Structure:}
\begin{itemize}
    \item The primal has one equality constraint. Therefore, the corresponding dual variable \( y_1 \) is unrestricted in sign.
    \item The second primal constraint is a \( \geq \) inequality, so its corresponding dual variable \( y_2 \) must be \( \geq 0 \).
    \item The third primal constraint is a \( \leq \) inequality, so its corresponding dual variable \( y_3 \) must be \( \leq 0 \).
    \item The dual constraints come from the coefficients of the primal variables. Since the primal is a maximization, the dual constraints are \( \geq \).
    \item The dual objective coefficients come from the right-hand sides of the primal constraints: \( 2, 3, 1 \).
\end{itemize}

\textbf{Summary:}  
This example demonstrates how mixed constraint types in the primal (\( =, \geq, \leq \)) lead to a dual with mixed variable sign restrictions. The primal variable bounds also affect whether dual constraints are inequalities or equalities. Understanding these mappings is crucial for correctly constructing dual problems.

\subsubsection*{Complicated Dual Derivation: A Second Look via Aggregation}
{}\label{example:complicated-dual}
Consider the following primal linear program:
\[
\begin{aligned}
\text{maximize} \quad & 2x_1 + x_2 \\
\text{subject to} \quad
& x_1 + x_2 = 2 \quad \text{(constraint 1)} \\
& 2x_1 - x_2 \geq 3 \quad \text{(constraint 2)} \\
& x_1 - x_2 \leq 1 \quad \text{(constraint 3)} \\
& x_1 \geq 0
\end{aligned}
\]

We now derive the dual using the \textbf{aggregation method}.

\vspace{0.5em}
\textbf{Step 1: Multiply constraints by dual variables.}

We introduce dual variables \( y_1, y_2, y_3 \) corresponding to constraints 1, 2, and 3, respectively.

We aggregate the constraints by forming a linear combination:
\[
y_1(x_1 + x_2 = 2) + y_2(2x_1 - x_2 \geq 3) + y_3(x_1 - x_2 \leq 1)
\]

Before combining, we must ensure all constraints are oriented as \( \leq \) inequalities, because we're maximizing in the primal. So:

- Equality constraint: no change needed, but equality implies \( y_1 \in \mathbb{R} \)
- \( \geq \): multiply by \( y_2 \geq 0 \)
- \( \leq \): multiply by \( y_3 \leq 0 \)

Now aggregate:
\[
\begin{aligned}
& y_1(x_1 + x_2) + y_2(2x_1 - x_2) + y_3(x_1 - x_2) \\
&= (y_1 + 2y_2 + y_3)x_1 + (y_1 - y_2 - y_3)x_2
\end{aligned}
\]

And the aggregated right-hand side is:
\[
2y_1 + 3y_2 + 1y_3
\]

So the aggregated constraint is:
\[
(y_1 + 2y_2 + y_3)x_1 + (y_1 - y_2 - y_3)x_2 \leq 2y_1 + 3y_2 + y_3
\]

\vspace{0.5em}
\textbf{Step 2: Ensure this constraint dominates the primal objective.}

To upper-bound the primal objective \( 2x_1 + x_2 \), we want:
\[
(y_1 + 2y_2 + y_3)x_1 + (y_1 - y_2 - y_3)x_2 \geq 2x_1 + x_2 \quad \text{for all feasible } x_1, x_2 \geq 0
\]

Since \( x_1, x_2 \geq 0 \), domination means:
\[
\begin{aligned}
y_1 + 2y_2 + y_3 &\geq 2 \quad \text{(coefficient on } x_1) \\
y_1 - y_2 - y_3 &\geq 1 \quad \text{(coefficient on } x_2)
\end{aligned}
\]

\vspace{0.5em}
\textbf{Step 3: Formulate the dual.}

The dual objective is the right-hand side of the aggregated constraint:
\[
\text{minimize } 2y_1 + 3y_2 + y_3
\]

With constraints:
\[
\begin{aligned}
y_1 + 2y_2 + y_3 &\geq 2 \\
y_1 - y_2 - y_3 &\geq 1 \\
y_2 &\geq 0 \quad \text{(from } \geq \text{ primal constraint)} \\
y_3 &\leq 0 \quad \text{(from } \leq \text{ primal constraint)} \\
y_1 &\in \mathbb{R} \quad \text{(from } = \text{ constraint)}
\end{aligned}
\]

\vspace{0.5em}
\textbf{Conclusion:}

This derivation shows how:
\begin{itemize}
  \item \emph{Constraint direction} in the primal determines the sign of the corresponding dual variable.
  \item \emph{Sign of a primal variable} (e.g., \( x_1 \geq 0 \)) determines the direction of the dual constraint: to ensure domination for non-negative variables, we require the dual combination to be \(\geq\) the primal coefficients.
  \item The \emph{dual objective} comes from aggregating the right-hand sides, weighted by the dual variables.
\end{itemize}
This example illustrates how careful sign and inequality management gives rise to a valid dual that bounds the primal objective from above.





\textbf{Question:} What happens if variable $x_1$ instead was constrained as $x_1 \leq 0$?



